\documentclass[12pt,a4paper,twoside,openright,titlepage,final]{article}
\usepackage{fontspec}
\usepackage{amsmath}
\usepackage{amsfonts}
\usepackage{amssymb}
\usepackage{makeidx}
\usepackage{graphicx}
\usepackage[hidelinks,unicode=true]{hyperref}
\usepackage[spanish,es-nodecimaldot,es-lcroman,es-tabla,es-noshorthands]{babel}
\usepackage[left=3cm,right=2cm, bottom=4cm]{geometry}
\usepackage{natbib}
\usepackage{microtype}
\usepackage{ifdraft}
\usepackage{verbatim}
\usepackage[table,xcdraw]{xcolor}
\usepackage[obeyDraft]{todonotes}
\ifdraft{
	\usepackage{draftwatermark}
	\SetWatermarkText{BORRADOR}
	\SetWatermarkScale{0.7}
	\SetWatermarkColor{red}
}{}
\usepackage{booktabs}
\usepackage{longtable}
\usepackage{calc}
\usepackage{array}
\usepackage{caption}
\usepackage{subfigure}
\usepackage{footnote}
\usepackage{url}
\usepackage{tikz}
\usepackage{pdflscape}
\usepackage{minted}
\usepackage{multirow}


%\setsansfont[Ligatures=TeX]{texgyreadventor}
%\setmainfont[Ligatures=TeX]{texgyrepagella}

%*******************************************************
%                 NO MODIFICAR
\newcommand*{\FSfont}[1]{%
  \fontencoding{T1}\fontfamily{#1}\selectfont}

\newlength{\tpheight}\setlength{\tpheight}{0.9\textheight}
\newlength{\txtheight}\setlength{\txtheight}{0.9\tpheight}
\newlength{\tpwidth}\setlength{\tpwidth}{0.9\textwidth}
\newlength{\txtwidth}\setlength{\txtwidth}{0.9\tpwidth}
\newlength{\drop}
%*******************************************************

% Crea una portada con los siguientes parámetros
%
% #1 : Título 
% #2 : Subtítulo
% #3 : Subsubtítulo
% #4 : Autor(es)
% #5 : Lugar
%

\newcommand*{\portada}[5]{
\begin{titlepage}
\begingroup
\vspace*{1cm}
\drop = 0.2\txtheight
\centering
\vfill
{\Huge \scshape #1}\\[\baselineskip]
{\Large \textbf{#2}}\\[\baselineskip]
{\Large \scshape #3}\\[\baselineskip]
\vspace*{0.3cm}
{\large \textit{#4}}\\[0.5\drop]
\includegraphics[scale=0.35]{./imagenes/logoURJC.jpg}
\vspace*{1.5cm}

{\large \scshape #5, \today} \par
\begin{center}
\end{center}
\vfill\null
\endgroup
\end{titlepage}
}
 %*****************************************************
 


\author{José Ignacio Escribano}

\setlength{\parindent}{0pt}

\begin{document}

\pagenumbering{alph}
\setcounter{page}{1}

\portada{Práctica III}{Simulación y Metaheurísticas}{Metaheurísticas}{José Ignacio Escribano}{Móstoles}

\listoftables
\thispagestyle{empty}
\newpage

\listoffigures
\thispagestyle{empty}
\newpage

\tableofcontents
\thispagestyle{empty}
\newpage


\pagenumbering{arabic}
\setcounter{page}{1}

\section{Introducción}

En esta práctica implementaremos una de las metaheurísticas planteadas para el problema del p-hubs.


\section{Resolución de la práctica}

A continuación mostraremos la implementación del recocido simulado para el problema del p-hub. El código se encuentra en el método estático \texttt{recocidoSimulado} de la clase \texttt{Práctica3}.

\begin{minted}[linenos=true, numberblanklines=true, breaklines=true]{java}
static Solución recocidoSimulado(Enfriamiento.TIPOS_ENFRIAMIENTO tipo, Solución sol_ini, InstanciaPHub instancia,
int temp_ini, int nrep, double alpha, long tiempo) {

    double tk = 0;

    Solución x = sol_ini;

    int i = 1;

    switch (tipo) {
        case BOLTZMANN: tk = Enfriamiento.crierioBoltzmann(temp_ini, alpha, i); break;
        case CAUCHY: tk = Enfriamiento.esquemaCauchy(temp_ini, alpha, i); break;
        case LUNDYyMEES: tk = Enfriamiento.esquemaCauchy(temp_ini, alpha, i); break;
        case DESCENSO_GEOMETRICO: tk = Enfriamiento.crierioBoltzmann(temp_ini, alpha, i);
    }

    long fin = System.currentTimeMillis() + tiempo;

    double delta;

    do {
        int m = 0;
        do {
            // Generamos número aleatorio entre 0 y número de soluciones de la vecindad - 1
            Random r = new Random();
            List<Solución> vecindad = PHub.generarVecindad(x, instancia);
            int n = r.nextInt(vecindad.size());

            Solución y = vecindad.get(n);

            delta = y.getObjetivo() - x.getObjetivo();

            if (delta < 0) {
                x = y;
            } else {
                double u = r.nextDouble();
                if (u <= Math.exp(-delta / tk)) {
                    x = y;
                }
            }
            m++;
        } while (m != nrep);
        
        switch (tipo) {
            case BOLTZMANN: tk = Enfriamiento.crierioBoltzmann(temp_ini, alpha, i); break;
            case CAUCHY: tk = Enfriamiento.esquemaCauchy(temp_ini, alpha, i); break;
            case LUNDYyMEES: tk = Enfriamiento.esquemaCauchy(temp_ini, alpha, i); break;
            case DESCENSO_GEOMETRICO: tk = Enfriamiento.crierioBoltzmann(temp_ini, alpha, i);
        }
        i++;
    } while (System.currentTimeMillis() <= fin);

    return x;
}
\end{minted}

Esta función recibe un tipo de enfriamiento, una solución inicial, una instancia, una temperatura inicial, un número de repeticiones, un valor (parámetro de la función de enfriamiento) y el tiempo de ejecución del programa, en milisegundos.\\

En la línea 6 se establece que \texttt{x} es la solución inicial, y desde la línea 10 a la 16 se establece qué función de enfriamiento se utilizará para el recocido simulado.\\

Se elige una solución al azar de la vecindad de \texttt{x}, que llamaremos \texttt{y}, y se calcula la diferencia entre la solución nueva y la inicial. A este valor lo llamemos \texttt{delta}. Si \texttt{delta} es menor que cero, es decir $\delta = f(y) - f(x) < 0 \iff f(y) < f(x)$, es decir, se mejora la solución actual y se guarda. En caso contrario (si no se mejora la solución) se genera un número aleatorio en el intervalo $(0,1)$. Si este número aleatorio es menor que $e^{-\delta/t_k}$, se guarda la solución \texttt{y}. Notar que $t_k$ es el valor de la temperatura en la iteración $k$. Este proceso se repite \texttt{nrep} veces, y todo lo anterior se repite durante el tiempo especificado. Éste es nuestro criterio de parada.\\

En cuanto la funciones de enfriamiento se ha definido una clase \texttt{Enfriamiento} dentro de la clase \texttt{Práctica3}. Se han implementado cuatro funciones de enfriamiento, que vienen dadas por las siguientes ecuaciones:

\begin{itemize}
\item Descenso geométrico: $t_{i+1} = \alpha t_i, \ \alpha \in [0.8, 0.99]$
\item Criterio de Boltzmann: $t_i = \dfrac{t_0}{1 + \log i}$
\item Esquema de Cauchy $t_i = \dfrac{t_0}{1+i}$
\item Lundy y Mees:  $t_{i+1} = \dfrac{t_i}{1 + \beta t_i}$, con $\beta$ muy pequeño
\end{itemize}

El código del criterio de Boltzmann se muestra a continuación:

\begin{minted}[linenos=true, numberblanklines=true, breaklines=true]{java}
public static double crierioBoltzmann(double t0, double alpha, int i) {
    return t0 / (1 + Math.log(i));
}
\end{minted}

Esta función recibe tres parámetros: la temperatura inicial, el valor del parámetro $\alpha$ y la iteración i. Se devuelve directamente el valor dado por la fórmula indicada anteriormente.\\ 

De forma similar se implementan los demás métodos.

\subsection{Análisis de los resultados}

Para evaluar este método ejecutamos durante 1 minuto (60 milisegundos) cada una de las instancias con cada uno de las funciones de enfriamiento.\\

Los parámetros usados han sido los siguientes:

\begin{itemize}
	\item Temperatura inicial = 10000
	\item $\alpha = \begin{cases}
	0.99, \text{ si es descenso geométrico}\\
	0.001, \text{ si es Lundi y Mees}
	\end{cases}$
	\item Número de repeticiones = 250
	\item Tiempo de ejecución (condición de parada) = 60000 (1 minuto por instancia)
\end{itemize}

Los resultados de la ejecución se pueden ver en la Tabla~\ref{tbl:resultados}. En amarillo se muestra el mejor valor de la función objetivo de cada instancia.\\

\begin{table}[htbp!]
	\centering
	\caption{Resultados de la ejecución del recocido simulado}
	\label{tbl:resultados}
	\begin{tabular}{@{}cccccc@{}}
		\toprule
		& \multicolumn{4}{c}{Función de enfriamiento}                                                                                                                                                                                                               &                         \\ \cmidrule(lr){2-5}
		\multirow{-2}{*}{Instancia} & \begin{tabular}[c]{@{}c@{}}Descenso \\ geométrico\end{tabular} & \begin{tabular}[c]{@{}c@{}}Criterio de\\ Boltzmann\end{tabular} & \begin{tabular}[c]{@{}c@{}}Esquema \\ de Cauchy\end{tabular} & \begin{tabular}[c]{@{}c@{}}Lundi \\ y Mees\end{tabular} & \multirow{-2}{*}{Media} \\ \midrule
		1                           & \cellcolor[HTML]{FFFE65}818.9541                               & 883.8607                                                        & 939.1282                                                     & 891.3631                                                & 883.3266                \\
		2                           & 1186.6325                                                      & 1438.0429                                                       & 1000.0160                                                    & \cellcolor[HTML]{FFFE65}974.8233                        & 1149.8787               \\
		3                           & 1152.2130                                                      & \cellcolor[HTML]{FFFE65}996.1984                                & 1103.6910                                                    & 1198.2770                                               & 1112.5949               \\
		4                           & 1026.4867                                                      & 1009.7156                                                       & 1042.9046                                                    & \cellcolor[HTML]{FFFE65}972.5211                        & 1012.9071               \\
		5                           & 1018.6530                                                      & 1033.4533                                                       & \cellcolor[HTML]{FFFE65}889.6963                             & 922.5584                                                & 966.0903                \\
		6                           & 1191.6062                                                      & 1049.2819                                                       & \cellcolor[HTML]{FFFE65}905.0196                             & 1179.4797                                               & 1081.3469               \\
		7                           & 1142.4889                                                      & 1140.0351                                                       & 1472.7656                                                    & \cellcolor[HTML]{FFFE65}914.2835                        & 1167.3933               \\
		8                           & 1077.5733                                                      & 1278.6855                                                       & 1149.6770                                                    & \cellcolor[HTML]{FFFE65}890.9676                        & 1099.2259               \\
		9                           & 1126.6379                                                      & \cellcolor[HTML]{FFFE65}920.1679                                & 971.2918                                                     & 922.4024                                                & 985.1250                \\
		10                          & 1158.27774                                                     & 1130.5921                                                       & 1121.0358                                                    & \cellcolor[HTML]{FFFE65}935.7108                        & 1086.4042               \\ \midrule
		Media                       & 1089.9524                                                      & 1088.0034                                                       & 1059.5226                                                    & 980.2387                                                & \multicolumn{1}{l}{}    \\ \bottomrule
	\end{tabular}
\end{table}

Se puede observar que en la primera instancia se obtiene mejor resultado con el método del descenso geométrico, en las instancias 3 y 9, el mejor método es el criterio de Boltzmann, en las instancias 5 y 6 se obtiene la mejor solución con el esquema de Cauchy, en las instancias 2, 4, 7,8 y 10 se obtiene la mejor solución con el método de Landi y Mees. Es decir, con el método del descenso geométrico se obtiene sólo 1 mejor instancia, con el criterio de Boltzmann y el esquema de Cauchy se obtienen 2 instancias y con el método de Landi y Mees se obtienen 5 instancias. Además, este método tiene la mejor media de tiempos. Parece que este método es el mejor de los cuatro (con los parámetros definidos anteriormente).\\

Necesitamos conocer cuál es mejor valor de $\beta$ que minimiza el valor de las instancias. Para ello, repetimos el proceso anterior para distintos parámetros de $\beta$, desde $1$ hasta $10^{-10}$. Los resultados se pueden ver en la Tabla


\section{Conclusiones}

\newpage


\end{document}